\subsection{Monocular Hand Pose and Shape Estimation}\label{subsec:sota_estimation}

Hand pose estimation from imagery is the core field this system inhabits.
The task is to recover, from pixel observations, a 3D description of hand
configuration --- joint positions, rotations, or a full parametric model ---
sufficient for gesture recognition, avatar animation, or motion logging.
Approaches differ principally in input modality, whether or not they commit
to a parametric model, and the compute required at inference time.

\subsubsection{Problem Formulation and Input Modalities}\label{subsubsec:sota_modalities}

Recovering 3D hand pose from 2D images is fundamentally ill-posed:
perspective projection discards depth, making the same image consistent with
infinitely many 3D configurations \cite{Panteleris2017UsingAS}.
The literature addresses this through three families of input, each resolving
the ambiguity differently.

\emph{RGB-D} methods augment the colour frame with a registered depth map
from a structured-light or time-of-flight sensor, making the 3D position of
each observed surface point directly measurable and reducing pose estimation
to a fitting problem \cite{Panteleris2017UsingAS}.
High accuracy is achievable without learned depth priors, but at the cost of
a dedicated sensor.

\emph{Multi-view RGB} resolves depth geometrically by triangulating
corresponding keypoints across two or more calibrated cameras.
Such rigs have been used to generate large-scale weakly-supervised training
annotations \cite{simon2017hand}, but the
synchronised, calibrated setup they require is unavailable to a general user.

\emph{Monocular RGB} is the most constrained setting: a single camera,
no depth channel, all 3D structure inferred from
appearance or temporal cues \cite{Panteleris2017UsingAS}.
Depth is ambiguous, self-occlusion is unavoidable, and scale varies with
camera distance.
This thesis adopts this modality because it imposes no hardware requirement
beyond the webcam a user already owns.

\subsubsection{Model-Free Keypoint Estimation}\label{subsubsec:sota_keypoint}

\emph{Model-free} detectors localise hand landmarks directly in image space
without fitting a parametric model at inference time, outputting a flat set
of 2D or 2.5D coordinates per hand.
Simon et al.\ \cite{simon2017hand} addressed the chronic scarcity of
annotated in-the-wild hand data---observing that no markerless hand
keypoint detectors for RGB images then existed---through \emph{multi-view
bootstrapping}: an initial detector is applied independently to every view
of a calibrated multi-camera rig, correct detections are robustly
triangulated in 3D and erroneous ones filtered as outliers, and the
reprojected triangulations serve as new labeled training data with which
an improved detector is iteratively retrained \cite{simon2017hand}.
A manual verification pass over the highest-scored frames is included,
though the authors note it typically requires only one to two minutes and
can largely be replaced by automatic heuristic filtering
\cite{simon2017hand}.
The resulting hand keypoint detector forms the hand-detection component
of OpenPose \cite{cao2021openpose}, and its 21-landmark topology is
adopted directly by MediaPipe Hands \cite{zhang2020mediapipehands}.
MediaPipe Hands, described in \autoref{subsec:mediapipe}, belongs to
this family.

The limitation shared across this family is that the output is a flat list
of image-space coordinates: no joint hierarchy, no canonical scale, no
rotation representation, and depth estimated only relative to the wrist
\cite{zhang2020mediapipehands}.
No downstream application requiring a skeletal pose or reusable motion
data can consume this output directly --- it is the representational gap
identified in \autoref{subsec:gap} that this thesis bridges.

\subsubsection{Model-Based Estimation: Direct MANO Regression}\label{subsubsec:sota_mano}

A second family closes this gap by predicting the parameters of a parametric
hand model directly from image features, bypassing the landmark stage.
Given its canonical status as a learned deformable hand prior
\cite{MANO:SIGGRAPHASIA:2017}, MANO is the output space adopted by most
modern regressors.

Boukhayma et al.\ \cite{boukhayma2019hand} presented the first end-to-end
deep learning method that predicts both 3D hand shape and pose from a single
unconstrained RGB image. A ResNet-50 encoder regresses MANO shape and pose
parameters; the decoder is a fixed, pre-computed MANO mesh generator coupled
with a weak-perspective reprojection module. The pipeline is trained end-to-end
with a combination of 3D joint supervision on controlled datasets and 2D joint
supervision on in-the-wild images, without requiring camera intrinsics at any
stage, achieving state-of-the-art 3D pose accuracy at the time of publication.

Hasson et al.\ \cite{hasson2019obman} extended the paradigm to joint
hand--object reconstruction by integrating MANO as a differentiable network
layer and simultaneously estimating object geometry with an AtlasNet decoder.
A novel contact loss---combining an attraction term that encourages hand-surface
contact with the object and a repulsion term that penalises
interpenetration---enforces physically plausible hand--object configurations
and improves grasp quality metrics over a baseline trained without contact
supervision.

FrankMocap \cite{rong2021frankmocap} incorporated a hand estimation module
into a whole-body system through a modular design in which independent
regression networks for the face, hands, and body are run separately and their
outputs are subsequently composed by an integration module. The hand module is
built on the hand component of the SMPL-X parametric model, uses a ResNet-50
encoder, and includes motion-blur data augmentation to improve robustness to the
blur artefacts common in full-body capture scenarios.

More recently, HaMeR \cite{pavlakos2024hamer} adopted a fully
transformer-based architecture pairing a large-scale Vision Transformer (ViT-H)
backbone with a transformer decoder that regresses MANO pose, shape, and camera
parameters. Trained on 2.7\,M examples consolidated from ten datasets containing
2D or 3D hand annotations, HaMeR consistently outperforms prior work on the
controlled FreiHAND and HO3D benchmarks and, by substantially larger margins, on
the authors' in-the-wild HInt benchmark. All methods in this family are evaluated
on GPU hardware in their respective publications.

The closest predecessor to the approach developed in this thesis is the system of
Zhou et al.\ \cite{zhou2020monocular}. It decomposes hand pose estimation into
two learned modules: \textit{DetNet}, a multi-task convolutional network with a
ResNet-50 backbone that predicts 2D and 3D joint positions from a single RGB
image by combining 2D-annotated in-the-wild data with 3D-annotated synthetic
data; and \textit{IKNet}, a seven-layer fully-connected neural network that maps
those 3D joint predictions to per-joint quaternion rotations in a single
feed-forward pass, enabling direct MANO mesh animation without iterative
optimisation. IKNet is trained jointly on motion-capture data from the MANO
dataset and on noisy 3D predictions produced by DetNet, so that it learns a
hand-pose prior while remaining robust to detection errors. Running on an NVIDIA
GTX\,1080\,Ti, the combined system achieves over 100\,fps, with DetNet
accounting for 8.9\,ms and IKNet for 0.9\,ms per frame.

\subsubsection{Real-Time Tracking Across the Hardware Spectrum}\label{subsubsec:sota_hardware}

The hardware alternatives to webcam-based tracking span a wide range of
accuracy and acquisition cost, making the accessibility constraint concrete.

The Leap Motion Controller tracks hand and finger position using three
infrared emitters and two infrared cameras. Controlled experiments
demonstrate axis-independent positional deviations below 0.2\,mm in static
conditions and an overall average accuracy of 0.7\,mm across all test
scenarios, broadly confirming the sub-millimetre accuracy the manufacturer
specifies \cite{weichert2013leap}. These results, however, require a
dedicated, purpose-built accessory placed at the interaction site.

Head-mounted platforms represent a second trajectory. The Oculus Quest
headset integrates on-device articulated hand tracking through four
egocentric monochrome fisheye cameras, enabling interaction without held
controllers or additional peripherals \cite{han2020megatrack}. The system
achieves real-time performance at 60\,Hz on a PC and 30\,Hz on a mobile
processor, yet this capability is inseparable from the headset itself
\cite{han2020megatrack}.

Each alternative delivers tracking fidelity that a monocular RGB approach
cannot match.
None operates on hardware a user already possesses.
The design constraint of this thesis --- a standard webcam and a
general-purpose CPU --- is therefore a deliberate commitment to the broadest
possible user base.
